% ============================================================
%  DCN Finals Reviewer
%  Polytechnic University of the Philippines
%  Data Communication and Networking | June 2026
% ============================================================
\documentclass[9pt]{article}

% -- Encoding & Fonts ----------------------------------------
\usepackage[T1]{fontenc}
\usepackage[utf8]{inputenc}
\usepackage{lmodern}

% -- Page geometry (8.5 x 13 in — Philippine Long / Foolscap)
\usepackage[
  paperwidth=8.5in,
  paperheight=13in,
  top=0.55in,
  bottom=0.55in,
  left=0.45in,
  right=0.45in,
  headheight=18.65pt,
  headsep=0.12in,
  footskip=0.3in
]{geometry}

% -- Colour --------------------------------------------------
\usepackage{xcolor}
\definecolor{themeColor}{HTML}{660100}
\definecolor{InkBlack}{HTML}{1A1A1A}
\definecolor{MutedGray}{HTML}{66615B}
\definecolor{SurfaceWhite}{HTML}{FFFFFF}
\definecolor{AnswerRed}{HTML}{8B1A1A}
\definecolor{LightRedBg}{HTML}{FDF2F2}
\color{InkBlack}

% -- Layout --------------------------------------------------
\usepackage{multicol}
\setlength{\columnsep}{0.28in}
\setlength{\columnseprule}{0.4pt}
\def\columnseprulecolor{\color{themeColor!30}}

% -- Section Headings ----------------------------------------
\usepackage{titlesec}
\titleformat{\section}
  {\large\bfseries\color{themeColor}}{\thesection.}{0.5em}{}[\vspace{-0.4em}\rule{\linewidth}{0.5pt}]
\titleformat{\subsection}
  {\normalsize\bfseries\color{themeColor}}{\thesubsection.}{0.5em}{}
\titlespacing*{\section}{0pt}{0.7em}{0.2em}
\titlespacing*{\subsection}{0pt}{0.5em}{0.15em}

% -- Lists ---------------------------------------------------
\usepackage{enumitem}
\setlist{nosep, leftmargin=1.2em}

% -- Boxes ---------------------------------------------------
\usepackage[most]{tcolorbox}

% -- Links ---------------------------------------------------
\usepackage[colorlinks=true, linkcolor=themeColor, urlcolor=themeColor]{hyperref}

% -- Header / Footer -----------------------------------------
\usepackage{fancyhdr}
\pagestyle{fancy}
\fancyhf{}
\lhead{\footnotesize\color{themeColor}\bfseries DCN Finals Reviewer}
\rhead{\footnotesize\color{MutedGray} Data Communication \& Networking}
\cfoot{\small\color{MutedGray}\thepage}
\renewcommand{\headrulewidth}{0.4pt}
\renewcommand{\headrule}{\hbox to\headwidth{\color{themeColor}\leaders\hrule height \headrulewidth\hfill}}

% -- Misc ----------------------------------------------------
\usepackage{microtype}
\usepackage{tikz}
\usepackage{eso-pic}
\raggedbottom
\raggedcolumns
\setlength{\parindent}{0pt}
\setlength{\parskip}{0pt}



% ============================================================
%  CUSTOM COMMANDS
% ============================================================

% Answer tag — red background, white bold text
\newcommand{\answer}[1]{%
  \par\nobreak\vspace{1pt}%
  \colorbox{AnswerRed}{\textcolor{white}{\bfseries\small\strut\ #1\ }}%
  \par\vspace{0pt}%
}

% Highlighted important word/phrase — bold + theme colour
\newcommand{\hl}[1]{\textbf{\textcolor{themeColor}{#1}}}

% Explanation box — subtle red-tinted background
\newcommand{\explain}[1]{%
  \par\nobreak\vspace{0pt}%
  \begin{tcolorbox}[
    colback=LightRedBg,
    colframe=themeColor!25,
    boxrule=0.3pt,
    arc=2pt,
    left=5pt, right=5pt,
    top=2pt, bottom=2pt,
    before skip=1pt,
    after skip=2pt
  ]
  {\small\textit{#1}}
  \end{tcolorbox}
}

% Question item
\newcounter{qnum}
\newcommand{\question}[1]{%
  \stepcounter{qnum}%
  \par\vspace{10pt}%
  {\bfseries\color{themeColor}\theqnum.}\ #1%
}

% ============================================================
\begin{document}
\thispagestyle{fancy}

% -- Title Block ---------------------------------------------
\begin{center}
  {\LARGE\bfseries\color{themeColor} DCN Finals Reviewer\par}
  \vspace{0.08cm}
  {\normalsize\itshape Data Communication and Networking\par}
  \vspace{0.04cm}
  {\small Polytechnic University of the Philippines\par}
  \vspace{0.06cm}
  {\color{themeColor}\rule{0.5\textwidth}{1pt}}
\end{center}
\vspace{0.15cm}

\begin{multicols}{2}

% ============================================================
\section{IPv6 Transition \& Address Types}
% ============================================================

\question{The layer that is closest to the \hl{end user}.}

\answer{Application Layer}

\explain{The Application layer is the topmost layer of both the OSI and TCP/IP models, providing the interface through which users interact with network services such as HTTP, FTP, and DNS.}

% --

\question{The devices run both \hl{IPv4} and \hl{IPv6} protocol stacks simultaneously.}

\answer{Dual Stack}

\explain{Dual stack allows a single device to process both IPv4 and IPv6 packets independently on the same interface.}

% --

\question{A method of transporting an \hl{IPv6 packet} over an \hl{IPv4 network}. The IPv6 packet is \hl{encapsulated} inside an IPv4 packet.}

\answer{Tunneling}

\explain{Tunneling wraps (encapsulates) an IPv6 packet inside an IPv4 header so it can traverse IPv4-only infrastructure without modification.}

% --

\question{\hl{Network Address Translation 64 (NAT64)} allows IPv6-enabled devices to communicate with IPv4-enabled devices using a \hl{translation technique} similar to NAT for IPv4.}

\answer{Translation}

\explain{Translation converts packet headers between IPv6 and IPv4 formats, enabling direct communication between hosts running different IP versions.}

% --

\question{Uniquely identifies an \hl{interface} on an IPv6-enabled device.}

\answer{Unicast}

\explain{A unicast address identifies exactly one interface; packets sent to a unicast address are delivered to that single destination.}

% --

\question{Used to send a single IPv6 packet to \hl{multiple destinations}.}

\answer{Multicast}

\explain{Multicast delivers a packet to all members of a specific group, reducing bandwidth usage compared to sending individual copies to each recipient.}

% --

\question{Any IPv6 unicast address that can be assigned to \hl{multiple devices}. A packet sent to this address is routed to the \hl{nearest device} having that address.}

\answer{Anycast}

\explain{Anycast uses routing metrics to forward the packet to the topologically closest interface sharing the same address, improving latency and redundancy.}

% --

\question{Similar to a \hl{public IPv4 address}. These are \hl{globally unique, internet-routable} addresses.}

\answer{Global Unicast Address (GUA)}

\explain{GUAs are assigned from the 2000::/3 range and are equivalent to public IPv4 addresses, allowing devices to be reached across the global Internet.}

% --

\question{Required for every \hl{IPv6-enabled device} and used to communicate with other devices on the same \hl{local link}. LLAs are \hl{not routable} and are confined to a \hl{single link}.}

\answer{Link-Local Address (LLA)}

\explain{Link-local addresses (fe80::/10) are automatically configured on every IPv6 interface and are used for neighbor discovery and routing protocol exchanges within a single network segment.}

% ============================================================
\section{IPv4 Addressing}
% ============================================================

\question{How many total bits make up a standard \hl{IPv4 address} structure?}

\answer{32}

\explain{An IPv4 address is a 32-bit number, typically written as four decimal octets separated by dots (e.g., 192.168.1.1).}

% --

\question{Which component is used by a network device to differentiate the \hl{network portion} from the \hl{host portion} of an IPv4 address?}

\answer{Subnet Mask}

\explain{The subnet mask is a 32-bit value where contiguous 1-bits mark the network portion and 0-bits mark the host portion of an address.}

% --

\question{What is the logical result of the \hl{ANDing process} when a binary 1 is compared to a binary 0?}

\answer{0}

\explain{The AND operation returns 1 only when both bits are 1; any other combination (1\,\&\,0, 0\,\&\,1, 0\,\&\,0) yields 0.}

% --

\question{What is the \hl{prefix length} for a subnet mask of \hl{255.255.255.0}?}

\answer{/24}

\explain{Three octets of all 1s (8 $\times$ 3) equals 24 bits set in the mask.}

% --

\question{According to \hl{RFC 1918}, which of the following is a \hl{private IPv4 address} range?}

\answer{192.168.0.0/16}

\explain{This range is reserved for internal network use and is not globally routable.}

% ============================================================
\section{Network Security Threats}
% ============================================================

\question{The \hl{discovery and mapping} of systems, services, or vulnerabilities.}

\answer{Reconnaissance Attacks}

\explain{Reconnaissance is the information-gathering phase where attackers probe networks using tools like ping sweeps and port scans to identify potential targets.}

% --

\question{The unauthorized \hl{manipulation of data}, system access, or \hl{user privileges}.}

\answer{Access Attacks}

\explain{Access attacks exploit authentication weaknesses or software vulnerabilities to gain unauthorized entry to systems or escalate privileges.}

% --

\question{The \hl{disabling or corruption} of networks, systems, or services.}

\answer{Denial of Service}

\explain{DoS attacks flood a target with traffic or exploit protocol weaknesses to exhaust resources, rendering the service unavailable to legitimate users.}

% ============================================================
\section{Malware Types}
% ============================================================

\question{A type of malware that propagates by \hl{inserting a copy of itself} into, and becoming part of, another program. It spreads from one computer to another, leaving infections as it travels.}

\answer{Viruses}

\explain{Viruses require a host file to propagate; they attach to legitimate programs and execute when the host program runs.}

% --

\question{Similar to viruses in that they \hl{replicate functional copies} of themselves and can cause the same type of damage. In contrast to viruses, which require the spreading of an infected host file, this is a \hl{standalone software} and does not require a \hl{host program} or human help to propagate.}

\answer{Worms}

\explain{Worms are self-propagating and exploit network vulnerabilities to spread automatically without user interaction or a host executable.}

% --

\question{A harmful piece of software that \hl{looks legitimate}. They do not reproduce by infecting other files. They \hl{self-replicate} and spread through \hl{user interaction} such as opening an email attachment or downloading and running a file from the internet.}

\answer{Trojan Horses}

\explain{Trojans disguise themselves as useful software to trick users into installing them, after which they can create backdoors, steal data, or download additional malware.}

% --

\question{Malware is short for\ldots}

\answer{Malicious Software}

\explain{``Malware'' is the umbrella term for any software intentionally designed to cause damage, including viruses, worms, trojans, ransomware, and spyware.}

% ============================================================
\section{Physical \& Infrastructure Threats}
% ============================================================

\question{Includes \hl{physical damage} to servers, routers, switches, cabling plant, and workstations.}

\answer{Hardware Threats}

\explain{Hardware threats encompass any event—accidental or deliberate—that physically damages networking equipment, such as vandalism, theft, or natural disasters.}

% --

\question{Includes \hl{temperature extremes} (too hot or too cold) or \hl{humidity extremes} (too wet or too dry).}

\answer{Environmental Threats}

\explain{Environmental threats affect equipment reliability; excessive heat causes overheating, while high humidity can lead to condensation and corrosion on circuit boards.}

% --

\question{Includes \hl{voltage spikes}, insufficient supply voltage (\hl{brownouts}), unconditioned power (\hl{noise}), and \hl{total power loss}.}

\answer{Electrical Threats}

\explain{Electrical threats are mitigated by using UPS (uninterruptible power supplies), surge protectors, and power conditioners to maintain clean, stable electricity.}

% --

\question{Includes \hl{poor handling} of key electrical components (\hl{electrostatic discharge}), lack of critical spare parts, poor cabling, and poor labeling.}

\answer{Maintenance Threats}

\explain{Maintenance threats arise from neglecting best practices such as ESD wrist straps, proper cable management, and keeping an adequate inventory of spare parts.}

% ============================================================
\section{Attack Types}
% ============================================================

\question{Implemented using \hl{brute force}, \hl{Trojan horse}, and \hl{packet sniffers}.}

\answer{Password Attacks}

\explain{Password attacks attempt to discover credentials through exhaustive guessing (brute force), keystroke logging (Trojans), or intercepting unencrypted credentials on the network (sniffers).}

% --

\question{A threat actor uses \hl{unauthorized privileges} to gain access to a system, possibly compromising the target.}

\answer{Trust Exploitation}

\explain{Trust exploitation abuses the trust relationships between systems—e.g., if System A trusts System B, compromising B gives the attacker access to A.}

% --

\question{A threat actor uses a \hl{compromised system} as a base for attacks against other targets. For example, using \hl{SSH (port 22)} to connect to a compromised host A. Host A is trusted by host B and, therefore, the threat actor can use \hl{Telnet (port 23)} to access it.}

\answer{Port Redirection}

\explain{Port redirection (pivoting) allows an attacker to tunnel traffic through a compromised host to reach otherwise inaccessible systems on the internal network.}

% --

\question{The threat actor is positioned \hl{in between two legitimate entities} in order to \hl{read or modify} the data that passes between the two parties.}

\answer{Man-in-the-Middle}

\explain{MitM attacks intercept communications in real time, enabling eavesdropping, session hijacking, or data manipulation without either party's knowledge.}

% ============================================================
\section{Firewall Filtering}
% ============================================================

\question{Prevents or allows access based on \hl{IP or MAC addresses}.}

\answer{Packet Filtering}

\explain{Packet filtering examines the source and destination addresses in each packet header and applies allow/deny rules without inspecting payload content.}

% --

\question{Prevents or allows access by specific \hl{application types} based on \hl{port numbers}.}

\answer{Application Filtering}

\explain{Application filtering controls traffic by matching well-known port numbers (e.g., port 80 for HTTP, port 443 for HTTPS) to allow or block specific services.}

% --

\question{Prevents or allows access to websites based on specific \hl{URLs or keywords}.}

\answer{URL Filtering}

\explain{URL filtering inspects the requested web address or embedded keywords and blocks access to sites that match a deny list or policy category.}

% --

\question{Incoming packets must be \hl{legitimate responses} to requests from internal hosts. \hl{Unsolicited packets} are blocked unless permitted specifically. SPI can also include the capability to recognize and filter out specific types of attacks, such as \hl{denial of service (DoS)}.}

\answer{Stateful Packet Inspection (SPI)}

\explain{SPI maintains a state table of active connections and only permits inbound traffic that corresponds to an established outbound session, providing stronger security than simple packet filtering.}

% ============================================================
\section{IPv6 Multicast \& Address Details}
% ============================================================

\question{A \hl{multicast group} that \hl{all IPv6-enabled devices} join. A packet sent to this group is received and processed by all IPv6 interfaces on the link or network.}

\answer{ff02::1 All-Nodes Multicast Group}

\explain{ff02::1 is the link-local scope all-nodes multicast address, analogous to a broadcast in IPv4—every IPv6 host on the link listens to this address.}

% --

\question{A \hl{multicast group} that \hl{all IPv6 routers} join. A router becomes a member of this group when it is enabled as an IPv6 router with the \texttt{ipv6 unicast-routing} global configuration command.}

\answer{ff02::2 All-Routers Multicast Group}

\explain{ff02::2 targets only routers on the local link; hosts send Router Solicitation messages to this address to discover available routers.}

% --

\question{What is the total address length of an IPv6 address and written in what format?}

\answer{128, Hexadecimal}

\explain{IPv6 utilizes a 128-bit address space to provide a significantly larger number of unique addresses compared to IPv4.}

% --

\question{The unofficial term used to refer to a segment of \hl{16 bits}, or \hl{four hexadecimal values}.}

\answer{Hextet}

\explain{A hextet refers to a four-hexadecimal digit segment representing 16 bits of the 128-bit address.}

% --

\question{According to \hl{Rule 1} of IPv6 compression, which zeros can be omitted?}

\answer{Leading Zeros in Any Hextet}

\explain{Only zeros at the start of a hextet can be omitted without making the address ambiguous.}

% --

\question{Which of the following is a valid compression of the address \texttt{2001:0db8:0000:1111:0000:0000:0000:0200}?}

\answer{2001:db8:0:1111::200}

\explain{This correctly applies Rule 1 (leading zeros) and Rule 2 (double colon for a contiguous string of all-zero hextets).}

% --

\question{True or False: IPv6 addresses are \hl{case-sensitive}.}

\answer{False}

\explain{IPv6 addresses are not case-sensitive and can be written in lowercase, uppercase, or a mix of both.}

% --

\question{Which address type did IPv6 \hl{eliminate} that was present in IPv4?}

\answer{Broadcast}

\explain{IPv6 does not have a broadcast address, using specific multicast groups instead to reduce unnecessary network noise.}

% --

\question{What is the recommended \hl{prefix length} for most IPv6 LANs?}

\answer{/64}

\explain{A 64-bit interface ID is required for SLAAC and simplifies network management and subnetting.}

% --

\question{\hl{Global Unicast Addresses (GUAs)} currently being assigned always begin with which \hl{binary bits}?}

\answer{001}

\explain{The binary prefix `001' corresponds to the 2000::/3 address space currently used for GUAs.}

% --

\question{The area between the Global Routing Prefix and the Interface ID. The \hl{Subnet ID} is used by an organization to \hl{identify subnets} within its site.}

\answer{Subnet ID}

\explain{This field allows organizations to create internal subnets within their allocated global routing prefix.}

% --

\question{The \hl{prefix}, or network, portion of the address that is assigned by the provider, such as an \hl{ISP}, to a customer or site. The global routing prefix will vary depending on ISP policies.}

\answer{Global Routing Prefix}

\explain{This portion forms the network identity that is routable on the global Internet.}

% --

\question{Equivalent to the \hl{host portion} of an IPv4 address. It is strongly recommended that in most cases /64 subnets should be used, which creates a \hl{64-bit} interface ID.}

\answer{Interface ID}

\explain{The Interface ID uniquely identifies a host on a specific subnet.}

% --

\question{What is the specific address range for \hl{IPv6 Link-Local Addresses (LLAs)}?}

\answer{fe80::/10}

\explain{LLAs are required for every interface and always fall within the fe80::/10 range.}

% --

\question{What prefix range is used by \hl{Unique Local Addresses (ULAs)}?}

\answer{fc00::/7 to fdff::/7}

\explain{ULAs are used for local addressing within a site and fall within this specific range.}

% --

\question{In the \hl{EUI-64 process}, what value is inserted into the \hl{middle of the 48-bit MAC address}?}

\answer{fffe}

\explain{The EUI-64 process takes the 48-bit MAC and inserts the 16-bit hex value `fffe' in the middle to create a 64-bit ID.}

% --

\question{Which \hl{ICMPv6 message} is sent by a router to provide \hl{addressing information} to hosts?}

\answer{Router Advertisement (RA)}

\explain{Routers send RA messages periodically or in response to an RS to provide prefix and gateway data.}

% --

\question{How often do IPv6 routers typically send out \hl{unsolicited ICMPv6 RA messages}?}

\answer{Every 200 seconds}

\explain{By default, IPv6-enabled routers send RA messages every 200 seconds to all nodes.}

\vspace{0.5em}
\begin{center}
  {\color{themeColor}\rule{0.6\linewidth}{0.8pt}}\\[4pt]
  {\footnotesize\itshape\color{MutedGray} End of Reviewer}
\end{center}

\end{multicols}
\end{document}
